\chapter{Specifications for Language-Model Test Generation}
\label{ch:testgen}

The previous chapter described how specifications are produced. This chapter
asks whether having them is worth anything, and answers the question against
the most demanding contemporary consumer of code-level information: the large
language model used to generate unit tests. The chapter is organised as a
single empirical study with two methodologically independent components. The
first establishes the \emph{fidelity} of the inferred specifications --- how
accurately they capture intended behaviour, measured against both human
judgment and the formal verifier --- and answers RQ1. The second establishes
their \emph{value} --- whether supplying them to a language model measurably
improves the tests it writes --- and answers RQ2. The two share a corpus but
are otherwise independent observations, and keeping them separate is essential
to the argument: the fidelity result tells us the specifications are good, and
the value result tells us they help, and neither subsumes the other.

\section{The Oracle Problem, Concretely}
\label{sec:tg-oracle}

A unit test has a stimulus and an oracle. The stimulus drives the code into
some state; the oracle asserts the state is correct. A language model generates
stimuli fluently --- given a method it readily produces inputs that exercise
its branches --- but the oracle requires a statement of \emph{intended}
behaviour that the code itself does not contain. Lacking such a statement, the
model does the only thing it can: it infers the oracle from the implementation,
and so asserts what the code does rather than what it should
do~\cite{konstantinou2024oracles}. When the implementation is correct this is
harmless; when it is buggy the model encodes the bug as the expected result,
producing a test that passes on broken code and will never catch the bug. Given
only a signature, with no body to infer from, the model falls back to vacuous
assertions --- that the method does not throw, that its result is not null ---
which compile and pass but distinguish almost nothing.

The hypothesis of this chapter is that an inferred specification, supplied to
the model as context, breaks this failure mode by carrying the missing
statement of intent. The specification says what the method should do
independently of what it does; a precondition delimits the valid inputs, a
postcondition states the expected output, and the model can assert against
these rather than against the implementation. The hypothesis is not obvious a
priori. The specification is itself derived from the code, so one might object
that it carries no information the model could not extract from the code
directly. The objection is answered by the experiment's design: the
specification is roughly an order of magnitude more compact than the body, it
states intent the body only implies, and --- as the results show --- it shifts
the model's output in a direction that source-code access does not fully
replicate at equivalent input size.

The objection deserves a fuller answer, because it touches the conceptual heart of
the study. It is true that the specification is derived from the code, and that in
an information-theoretic sense it contains nothing the code does not. But the model
is not an information-theoretic ideal; it is a bounded reasoner that attends to
what is salient in its input, and the same fact can be salient in one
representation and buried in another. A postcondition stating that a method returns
the maximum of its arguments makes the intent salient in five words; the same
intent is present in the method body but buried in a ternary expression the model
must recognise, interpret, and generalise into a statement of intent before it can
test against it. The specification does the recognising, interpreting, and
generalising in advance, and presents the result in the form the model can act on
directly. The value is not new information in the Shannon sense but \emph{distilled}
information --- the same intent, extracted from the implementation and stated
explicitly --- and the distillation is exactly what a model lacking a reliable
specification-production capability cannot do for itself, as the producer-side
comparison confirms. The study measures the value of the distillation, and that the
distillation helps even though it adds no information the code lacked is the precise
sense in which a specification is more than a summary of the code: it is the code's
intent made salient.

\section{Experimental Design}
\label{sec:tg-design}

The experiment holds the model fixed and varies only its input, isolating the
contribution of specifications. The same 312-method corpus is used throughout,
and the model is queried under four conditions.

\begin{description}
  \item[P1 --- Signature only.] The model receives the method signature alone.
        This is the baseline: the floor of what a model can do with no
        behavioural information beyond types.
  \item[P2 --- Signature with guidance.] The model receives the signature
        together with explicit instruction to cover edge cases. This condition
        is the crucial control, included to test the alternative hypothesis
        that any prompt enrichment, not specifications in particular, would
        account for an effect.
  \item[P3 --- Signature with specification.] The model receives the signature
        together with the inferred JML specification. This is the treatment
        condition.
  \item[P4 --- Full source.] The model receives the complete method body. This
        is the upper bound: with the body in hand the model has the test answer
        directly, so P4 bounds the achievable quality from above and is not a
        realistic deployment condition.
\end{description}

The model is Google's Gemini~2.0, queried at temperature 0.3 to balance
diversity against consistency, with five independent runs per method per
condition to control for sampling non-determinism. Prompts are kept minimal and
differ across conditions only in the inclusion of the specification (P3) or the
source (P4) or the guidance instruction (P2); no other prompt engineering is
applied, so that any difference in output is attributable to the input
variation rather than to prompt craft. The choice of a single model family is a
deliberate budget decision --- the experiment requires several thousand
generations and was repeated during pilot tuning, totalling on the order of
twenty-five thousand model invocations --- and the threats-to-validity
discussion (Section~\ref{sec:tg-threats}) addresses the single-model limitation
directly. The defended contribution is the \emph{ordering} of the conditions,
not the absolute scores, and the ordering is expected to be robust across
models on the evidence of multi-model studies in the
literature~\cite{schafer2023adaptive, yuan2024chattester, molinelli2025oracles}.

\subsection{Subject Selection}

The corpus is drawn from Apache Commons Lang 3.14, chosen for its diversity of
method types, its thorough Javadoc documentation (which enables the manual
fidelity validation), its maturity, its deterministic pure functions (well
suited to mutation testing), and its precedent in the research
literature~\cite{pacheco2007randoop, fraser2011evosuite}. Eleven classes were
selected to cover five of the six method categories of
Section~\ref{sec:inf-cat}: \texttt{NumberUtils}, \texttt{BooleanUtils}, and
\texttt{CharUtils} (utility); \texttt{MutableInt}, \texttt{MutableDouble}, and
\texttt{MutableBoolean} (state modification); \texttt{Pair} and
\texttt{MutablePair} (factory/delegate); \texttt{Validate} and \texttt{Range}
(control flow); and the \texttt{Mutable} interface. All public, package-private,
and protected methods of these classes were included, with no method excluded
after the fact, yielding 312 methods. The distribution across categories is
shown in Table~\ref{tab:tg-distribution}. The \emph{other} category is empty,
because Commons Lang contains no event handlers or callbacks --- a structural
feature of the corpus whose consequence for external validity is taken up in
Section~\ref{sec:tg-threats}.

\begin{table}[ht]
\centering
\caption{Distribution of the 312-method corpus by category. The \emph{other}
category contains zero methods because Apache Commons Lang includes no
event-handler or callback methods.}
\label{tab:tg-distribution}
\begin{tabular}{lrr}
\toprule
\textbf{Category} & \textbf{Count} & \textbf{\%} \\
\midrule
Utility Methods & 89 & 28.5\% \\
State Modification & 78 & 25.0\% \\
Accessors \& Mutators & 62 & 19.9\% \\
Control Flow & 48 & 15.4\% \\
Factory \& Delegate & 35 & 11.2\% \\
Other (callbacks, event handlers) & 0 & 0.0\% \\
\midrule
\textbf{Total} & \textbf{312} & \textbf{100\%} \\
\bottomrule
\end{tabular}
\end{table}

\subsection{Metrics}

Test quality is measured by four quantities. \emph{Test count} is the number of
test methods generated, averaged over the five runs. \emph{Compilation rate} is
the fraction of generated tests that compile. \emph{Pass rate} is the fraction
of compiling tests that pass against the unmutated code. \emph{Mutation score},
the headline metric, is the fraction of seeded mutants the suite kills, computed
with PIT~\cite{coles2016pitest} using its default operator set and a per-test
timeout of ten seconds. Equivalent mutants are filtered by PIT's default
heuristics, with the residual rate of trivially equivalent mutants estimated at
under five percent via the sampling procedure of Papadakis et
al.~\cite{papadakis2019mutation}. The mutation score is the metric of record
because it measures fault-detection power directly and penalises exactly the
oracle-free test that covers a line without checking its behaviour; test count
is reported alongside as a complementary metric, and the relationship between
the two is itself a finding.

\subsection{Statistical Analysis}

For each metric the analysis reports mean, standard deviation, median,
interquartile range, and 95\% bootstrap confidence intervals over ten thousand
iterations. Phase comparisons use paired $t$-tests, with Wilcoxon signed-rank
tests as a non-parametric corroboration where within-phase distributions depart
from normality, and Cohen's $d$ as the effect-size measure. Because the design
yields multiple pairwise contrasts --- six pairs of conditions across four
metrics, twenty-four tests in total --- a Bonferroni correction is applied at
family-wise $\alpha = 0.05$, giving a per-test threshold of approximately
0.0021. Reported significant differences survive this corrected threshold. A
power analysis based on a thirty-method pilot establishes that the design
detects an effect of 2.2 percentage points at 80\% power under the corrected
threshold; the observed headline effect of 40.7 percentage points yields a
post-hoc power exceeding 0.999.

\subsection{The Experimental Infrastructure}

The experiment is mediated by a test-generation harness that warrants description,
because its design determines what the measurement can and cannot capture. The
harness extracts each method from the corpus, constructs the prompt for the
condition under test --- signature alone, signature with guidance, signature with
specification, or full source --- queries the model, parses the generated test
code from the response, writes it into a test project, and records the outcome. The
parsing step is more delicate than it appears: a model's response wraps the test
code in prose and markdown fences, and the harness must extract the code robustly
across the variations in how the model formats its output. A parser too strict
discards valid code that happens to be formatted unexpectedly; a parser too lax
admits prose as code and fails to compile.

The harness's handling of the model's non-determinism is to run each condition five
times per method and record each run independently, so that the per-method outcome
is a distribution rather than a point. The downstream analysis aggregates these
distributions, and the separation of generation from analysis --- the harness
produces logs, the analysis scripts consume them --- is what makes the analysis
reproducible even though the generation is not. The harness also copies the
generated tests into a project configured for mutation testing, so that the same
artefacts the harness produced are the artefacts PIT measures, with no manual
intervention between generation and measurement that could introduce inconsistency.

The infrastructure's design reflects a lesson that recurs in the distribution work:
the steps a harness performs silently are the steps most likely to corrupt a
measurement. The code-extraction step, in particular, can silently discard valid
tests if the model's formatting drifts, depressing the test count for reasons
unrelated to the condition under test. The harness logs the raw response alongside
the extracted code precisely so that an unexpectedly low test count can be traced to
either a genuine failure of the model to generate tests or an extraction failure on
tests it did generate --- a distinction that matters for interpreting the results,
and one the downstream probe of Chapter~\ref{ch:compositional} would later need
when its extraction encountered exactly this kind of drift.

\subsection{Reproducibility and the Replication Package}

An experiment with several thousand stochastic model invocations raises an
obvious reproducibility concern, and the design addresses it on three fronts.
First, every model invocation is logged with its prompt, its raw response, and
the parsed test code, so that the path from input to measured outcome is
auditable for any of the thousands of generations. Second, the analysis is driven
by scripts that take the logged outputs as input and produce the reported tables
and statistics deterministically, so that the step from raw generations to
headline numbers is itself reproducible even though the generations are not.
Third, the replication package parameterises the model client, so that the
experiment can be re-run against a different model by changing a configuration
value, and the analysis scripts accept a model identifier as an input column,
which is the structure a multi-model replication would use.

The non-determinism that remains --- the model's sampling at temperature 0.3 ---
is handled statistically rather than eliminated. Five runs per method per
condition give a per-method mean whose standard error is small relative to the
between-condition effect, and the bootstrap confidence intervals quantify the
residual uncertainty without assuming a parametric form for the sampling
distribution. The pilot study established that increasing from five to ten runs
reduced the standard error by under four percent, which did not justify doubling
the model cost; five runs is the point at which additional sampling buys
negligible precision. The reproducibility claim is therefore not that a re-run
would produce identical tests --- it would not, the model being stochastic --- but
that a re-run would produce the same statistical conclusions, and that the path
from any given set of generations to those conclusions is fully scripted and
auditable.

\section{Fidelity of the Inferred Specifications (RQ1)}
\label{sec:tg-fidelity}

Before measuring whether specifications help, the study establishes that the
inferred specifications are accurate, under two independent oracles.

\subsection{Manual Validation Against Documented Intent}

Two authors independently validated the inferred specifications for all 312
methods against the Javadoc documentation, the source code, and the standard-
library specifications, labelling each clause as correct, incomplete,
incorrect, or overconstrained, with disagreements resolved by discussion.
Inter-rater reliability was high --- Cohen's $\kappa = 0.91$ for preconditions
and $0.87$ for postconditions --- which gives confidence that the labels
reflect a stable judgment rather than individual idiosyncrasy. The results
appear in Table~\ref{tab:tg-validation}.

\begin{table}[ht]
\centering
\caption{Manual validation of inferred specifications against documented
intent, over all 312 methods.}
\label{tab:tg-validation}
\begin{tabular}{lrrr}
\toprule
\textbf{Classification} & \textbf{Precond.} & \textbf{Postcond.} & \textbf{Overall} \\
\midrule
Correct & 294 (94.2\%) & 273 (87.6\%) & 567 (90.9\%) \\
Incomplete & 13 (4.2\%) & 29 (9.3\%) & 42 (6.7\%) \\
Incorrect & 4 (1.3\%) & 7 (2.2\%) & 11 (1.8\%) \\
Overconstrained & 1 (0.3\%) & 3 (1.0\%) & 4 (0.6\%) \\
\midrule
\textbf{Total} & 312 & 312 & 624 \\
\bottomrule
\end{tabular}
\end{table}

Precision is 94.2\% for preconditions and 87.6\% for postconditions. The lower
postcondition precision is attributable chiefly to incompleteness for methods
with complex return logic, where the engine captures part but not all of the
returned relationship. Recall, measured against documented behaviour, is 89.3\%
for preconditions and 78.1\% for postconditions; the lower postcondition recall
reflects the difficulty of inferring relationships involving collection
properties or algorithmic correctness, which the local syntactic patterns do
not reach. The asymmetry between precondition and postcondition accuracy is a
consistent theme: preconditions arise from explicit guards that the engine
recognises reliably, whereas postconditions must capture the full effect of a
method's logic, which is harder to do from syntax alone.

\subsection{The Manual Validation Protocol}

The manual validation deserves a fuller account, because the precision and recall
figures rest on it and a reader should know how the labels were produced. Two
authors independently examined the inferred specification for each of the 312
methods, with three reference materials in view: the method's Javadoc
documentation, which states the documented intent; the method's source, which
states the actual behaviour; and the standard-library specifications, which fix
the behaviour of any library methods the body calls. Each clause was assigned one
of four labels. \emph{Correct} meant the clause agrees with the documented intent
and the implementation. \emph{Incomplete} meant the clause is correct as far as it
goes but omits part of the documented behaviour --- the common case for
postconditions of methods with complex return logic. \emph{Incorrect} meant the
clause asserts something false. \emph{Overconstrained} meant the clause is true but
stronger than warranted --- a precondition forbidding inputs the method actually
handles.

The four-way labelling is more informative than a binary correct/incorrect split,
because it separates the two ways a clause can fall short of perfect: by claiming
too little (incomplete) or by claiming too much (incorrect or overconstrained). The
distinction maps onto the precision/recall decomposition: precision is the fraction
of emitted clauses that are correct, penalised by the incorrect and overconstrained
labels, while recall is the fraction of documented behaviour captured, penalised by
the incomplete label and by behaviours for which no clause was emitted at all. The
two failure modes have different consequences downstream --- an incorrect clause
misleads a consumer, while an incomplete one merely leaves value on the table ---
which is why the thesis reports them separately rather than collapsing them.

Disagreements between the two raters were resolved by discussion, and the
inter-rater reliability before resolution --- Cohen's $\kappa = 0.91$ for
preconditions and $0.87$ for postconditions --- is reported as evidence that the
labelling was reproducible rather than idiosyncratic. The slightly lower agreement
on postconditions reflects their greater subtlety: whether a postcondition is
``complete'' relative to documented intent is a judgment call on which careful
raters can differ, while whether a precondition correctly negates a guard is more
nearly mechanical. The high agreement on both, and the resolution of the residual
disagreements by discussion, is what licenses treating the manual labels as a
ground truth against which the formal discharge can be compared as an independent
oracle.

\subsection{Formal Discharge Against the Implementation}

Manual validation establishes agreement with documented intent; OpenJML's
Extended Static Checker establishes formal consistency with the implementation.
The two oracles are independent, as Section~\ref{sec:inf-validation} argued, and
a clause may satisfy either without the other. Table~\ref{tab:tg-discharge}
reports the discharge outcome over all inferred clauses.

\begin{table}[ht]
\centering
\caption{OpenJML ESC discharge over all inferred clauses. \emph{Discharged}
clauses are verified against the implementation; \emph{flagged} clauses fail
discharge and are surfaced for review; \emph{unsupported} clauses involve JML
constructs the verifier cannot yet handle.}
\label{tab:tg-discharge}
\begin{tabular}{lrrrr}
\toprule
\textbf{Clause type} & \textbf{Inferred} & \textbf{Discharged} & \textbf{Flagged} & \textbf{Unsupported} \\
\midrule
Preconditions     & 312 & 296 (94.9\%) & 11 (3.5\%) &  5 (1.6\%) \\
Postconditions    & 312 & 268 (85.9\%) & 22 (7.1\%) & 22 (7.1\%) \\
Loop invariants   &  53 &  43 (81.1\%) &  4 (7.5\%) &  6 (11.3\%) \\
Assignable        & 308 & 308 (100.0\%) &  0 (0.0\%) &  0 (0.0\%) \\
\midrule
\textbf{Overall}  & \textbf{985} & \textbf{915 (92.9\%)} & \textbf{37 (3.8\%)} & \textbf{33 (3.4\%)} \\
\bottomrule
\end{tabular}
\end{table}

The discharge rate exceeds the manual-validation precision in every category,
which is initially surprising but follows from the independence of the oracles:
the verifier accepts clauses the human raters judged incomplete, because its
oracle is consistency with the implementation, not completeness against the
documentation. The 37 flagged clauses --- chiefly subtle overflow-boundary
postconditions on signed-integer methods, plus one loop invariant where the
heuristic produced a non-decreasing rather than a monotonically decreasing
variant --- would have been false positives had they been retained silently;
they are routed to a review queue and excluded from the specifications supplied
to the model in P3. The 33 unsupported clauses are dominated by nested-quantifier
postconditions involving \texttt{\textbackslash sum} under conditional path
conditions; they are recorded as not formally verifiable but retained for the
P3 prompts because their content is syntactically valid JML and useful as oracle
information to the model, even though it is outside the verified-sound surface.
The residual non-discharge, examined in detail in Section~\ref{sec:tg-discussion},
falls into three categorically identifiable tractability limits of the SMT
solver rather than into arbitrary failures of the engine.

\subsubsection{Why Discharge Exceeds Manual Precision}

The discharge rate exceeding the manual-validation precision in every category is
the single most informative relationship between the two oracles, and it deserves
a careful reading because it would be paradoxical if the oracles measured the same
thing. They do not. Manual precision asks whether a clause matches the documented
\emph{intent}; discharge asks whether a clause is consistent with the
\emph{implementation}. A clause can be consistent with the implementation while
falling short of the documented intent --- this is exactly the
\emph{incomplete} label of the manual validation, a clause that is true but says
less than the documentation promises. Such a clause depresses the manual precision
(it is not fully \emph{correct} against intent) while passing discharge (it is
\emph{consistent} with the code), and the gap between the two rates is, to a first
approximation, the population of such incomplete-but-consistent clauses.

The relationship has a practical consequence for how the two oracles should be used
together. Discharge cannot replace manual validation, because a clause that
discharges may still be incomplete relative to intent, and a consumer relying on
completeness --- a verifier proving a strong property, a test generator seeking
thorough coverage --- needs the manual oracle's completeness judgment that
discharge does not provide. Manual validation cannot replace discharge, because a
clause the raters judge correct against intent may still conflict with the
implementation's actual arithmetic or aliasing in a corner the documentation does
not mention --- the subtle overflow boundary that discharge catches and human
review misses. The two are genuinely independent, catching disjoint classes of
error, and the thesis runs both precisely because neither subsumes the other. The
discharge-exceeds-precision relationship is the empirical proof of the independence:
were the oracles measuring the same property, their rates would track, and the
systematic gap shows they do not.

\subsection{The Comparison Baselines in Detail}

The comparison of the engine against Jdoctor and against direct language-model
inference (reported below) deserves a fuller account of how the baselines were
constructed, because a comparison is only as informative as the fairness of its
baselines. Jdoctor was run on the same 312-method corpus in its standard
configuration, extracting exception preconditions from the methods' Javadoc. Its
coverage is bounded by what the documentation states: a method whose Javadoc
documents a thrown exception yields a precondition, and a method whose
documentation is silent yields nothing, which is why Jdoctor produces
specifications for 45.2\% of the corpus against the engine's 99.0\%. The
comparison is not a competition the engine wins outright, because Jdoctor recovers
documented exceptional behaviour the engine misses --- 25 exception preconditions
encoded in natural language that no structural pattern reaches. The honest reading
is complementarity: the two tools draw on different sources, documentation and
code structure, and a combined workflow captures the union.

The language-model inference baseline was constructed by prompting the same model
used in the test-generation experiment to emit JML directly from each method's
source, under the same five-run protocol. The construction is deliberately
favourable to the model: it sees the full source, the richest possible input, and
is asked only to produce the specification. Even so, its precision falls short of
the engine's --- 81.3\% against 94.2\% for preconditions --- because it
hallucinates clauses that are plausible against the method name but false against
the body, and its run-to-run consistency is 72.4\% against the static engine's
100\%, because its sampling is stochastic. The comparison isolates the value of
grounding: the model and the engine see the same code, but the engine's output is
determined by the code's structure while the model's is a prediction, and the
prediction is less precise and less stable. This is the empirical basis for the
thesis's producer-consumer asymmetry, developed in the synthesis of
Chapter~\ref{ch:discussion}: the model is the weaker producer here and the
stronger consumer in the main experiment, and the two findings together motivate
the architecture that uses each where it is strong.

\subsection{Comparison With Other Inference Approaches}

Table~\ref{tab:tg-comparison} compares the engine with Jdoctor and with direct
language-model inference (the same Gemini model prompted to emit JML from
source).

\begin{table}[ht]
\centering
\caption{Comparison with Jdoctor and language-model specification inference.}
\label{tab:tg-comparison}
\begin{tabular}{lrrr}
\toprule
\textbf{Metric} & \textbf{JML-Inferrer} & \textbf{Jdoctor} & \textbf{LLM (Gemini 2.0)} \\
\midrule
Methods with specs & 309 (99.0\%) & 141 (45.2\%) & 312 (100\%) \\
Precision (precond.) & 94.2\% & 92\% & 81.3\% \\
Recall (precond.) & 89.3\% & 83\% & 88.7\% \\
Precision (postcond.) & 87.6\% & --- & 74.8\% \\
Recall (postcond.) & 78.1\% & --- & 82.1\% \\
Consistency (5 runs) & 100\% & 100\% & 72.4\% \\
\bottomrule
\end{tabular}
\end{table}

The engine produces non-trivial specifications for 99.0\% of the corpus, against
Jdoctor's 45.2\% --- Jdoctor requires documented exceptions and is silent where
they are absent. The two are complementary: Jdoctor captures 25 exception
preconditions the engine misses, typically encoded in natural language, while
the engine captures 37 the Jdoctor misses, undocumented in Javadoc. Against
direct language-model inference, the engine attains markedly higher precision
(94.2\% versus 81.3\%) and perfect run-to-run consistency (100\% versus 72.4\%),
because its proposals are grounded in concrete code structure whereas the model
may hallucinate plausible-but-false clauses and varies across runs. The
comparison is the empirical basis for the methodological caution stated in
Chapter~\ref{ch:background}: a language model is a strong consumer of grounded
specifications and a weaker producer of them.

\section{Value of Specifications for Test Generation (RQ2)}
\label{sec:tg-value}

With the fidelity of the specifications established, the study turns to whether
supplying them improves test generation. Table~\ref{tab:tg-results} reports the
per-class test counts and mutation scores for each condition, with overall
quality metrics in the lower panel.

\begin{table}[t]
\centering
\caption{Test-generation results by class: per-method average test counts over
five runs, and mutation scores, for each condition. The lower panel reports
overall quality metrics as the unweighted mean across all 312 methods.}
\label{tab:tg-results}
\small
\begin{tabular}{l|cc|cc|cc|cc}
\toprule
\multirow{2}{*}{\textbf{Class}} &
\multicolumn{2}{c|}{\textbf{P1: Sig}} &
\multicolumn{2}{c|}{\textbf{P2: Sig+Guide}} &
\multicolumn{2}{c|}{\textbf{P3: Sig+Spec}} &
\multicolumn{2}{c}{\textbf{P4: Source}} \\
\cline{2-9}
 & Tests & Mut.\% & Tests & Mut.\% & Tests & Mut.\% & Tests & Mut.\% \\
\midrule
NumberUtils    & 12 & 28 & 48 & 78 & 41 & 68 & 65 & 92 \\
BooleanUtils   & 10 & 31 & 45 & 82 & 38 & 71 & 62 & 95 \\
CharUtils      &  9 & 30 & 41 & 80 & 35 & 70 & 58 & 93 \\
MutableInt     & 10 & 26 & 40 & 89 & 36 & 69 & 59 & 100 \\
MutableDouble  &  8 & 23 & 35 & 85 & 32 & 65 & 52 & 97 \\
MutableBoolean &  7 & 25 & 33 & 84 & 30 & 66 & 50 & 96 \\
Pair           &  8 & 29 & 36 & 83 & 33 & 70 & 53 & 95 \\
MutablePair    &  8 & 27 & 34 & 82 & 31 & 68 & 51 & 94 \\
Validate       &  8 & 35 & 32 & 81 & 29 & 74 & 48 & 94 \\
Range          &  6 & 22 & 28 & 76 & 25 & 62 & 42 & 91 \\
\midrule
\textbf{Mean (per method)} & 9.0 & 27.5 & 38.0 & 81.8 & 33.5 & 68.2 & 54.7 & 94.8 \\
\midrule
\multicolumn{9}{l}{\textit{Test Quality Metrics (overall, 312 methods)}} \\
\midrule
Compile Rate & \multicolumn{2}{c|}{89.2\%} & \multicolumn{2}{c|}{87.8\%} & \multicolumn{2}{c|}{91.5\%} & \multicolumn{2}{c}{93.2\%} \\
Pass Rate & \multicolumn{2}{c|}{94.1\%} & \multicolumn{2}{c|}{92.3\%} & \multicolumn{2}{c|}{96.8\%} & \multicolumn{2}{c}{97.1\%} \\
Valid Tests & \multicolumn{2}{c|}{83.9\%} & \multicolumn{2}{c|}{81.1\%} & \multicolumn{2}{c|}{88.6\%} & \multicolumn{2}{c}{90.5\%} \\
\bottomrule
\end{tabular}
\end{table}

\subsection{The Principal Finding}

Specification-guided generation (P3) produces 272\% more tests than
signature-only generation (P1) --- 95\% bootstrap confidence interval
$[245, 299]$, paired $t$-test $p < 0.001$, Cohen's $d = 2.41$, confidence
interval $[2.08, 2.79]$ --- and achieves a mutation score 40.7 percentage
points higher: 68.2\% against 27.5\%, paired $t$-test $p < 0.001$, Cohen's
$d = 3.07$, confidence interval $[2.71, 3.46]$. The effect is large by any
conventional standard and survives the Bonferroni correction comfortably. Full
source access (P4) attains the highest absolute mutation score, 94.8\%, as
expected when the model is handed the test answer directly, but the comparison
that matters is not P3 against P4 in absolute terms; it is the comparison in
\emph{input density}. A JML contract typically occupies five to twelve lines of
annotation, against the thirty to one hundred fifty lines of body that P4
supplies. Within that order-of-magnitude smaller input, P3 captures roughly
sixty percent of the achievable gain --- $(68.2 - 27.5)/(94.8 - 27.5)$ --- and
the residual gap to P4 is largely accounted for by the body containing the
answer outright. P3 also attains higher compilation and pass rates than P1 or
P2, evidence that specifications constrain the model toward syntactically and
semantically valid test code rather than merely toward more of it.

The input-density argument deserves to be stated as a principle, because it is the
form in which the result generalises beyond the particular numbers. The value of an
input to a test generator is not its size but its information content about
intended behaviour per unit of size. A signature is small and carries little such
information; a body is large and carries much, but most of it is implementation
detail rather than intent, and crucially it carries the test answer, which makes it
an upper bound rather than a realistic input. A specification is small and carries
information that is almost entirely about intent --- it is the distillation of the
body's behaviour into the contract the body honours. The sixty percent of the
achievable gain that P3 captures within an input an order of magnitude smaller than
P4's is the quantitative expression of this density: the specification is the
information-dense input, and its density is why it captures most of the benefit at
a fraction of the size. This reframes the comparison with the source baseline: P4
wins on absolute score because it contains the answer, but it does so with an input
that no deployment would use, and the practical comparison is between the signature
a deployment always has and the specification a deployment can cheaply infer, where
the specification's density is decisive.

\subsection{Why More Tests Are Not Better Tests}

The most instructive comparison in the study is P2 against P3. The guidance
condition P2 produces the \emph{highest} per-method test count (38.0, against
P3's 33.5) yet a mutation score 13.6 percentage points \emph{lower} (68.2\%
against 81.8\% --- the relationship inverts when measured by mutation rather
than count). This is the result that the guidance control was designed to
surface, and it refutes the alternative hypothesis that any prompt enrichment
would explain the P3 effect: enrichment with edge-case guidance produces more
tests but weaker ones. Inspection of representative P2 outputs identifies three
mechanisms. First, the edge-case instruction prompts the model to enumerate
near-duplicates --- separate cases for \texttt{0}, \texttt{-0},
\texttt{Integer.MIN\_VALUE}, and \texttt{Integer.MIN\_VALUE + 1} --- that
inflate the count without expanding the set of killed mutants. Second, in the
absence of an oracle, edge-case tests default to weak assertions
(\texttt{assertNotNull}, \texttt{assertDoesNotThrow}) that compile and pass but
fail to distinguish mutants that alter return values. Third, P3 tests are
concentrated on clauses --- a test or two per inferred precondition and
postcondition --- producing higher mutant-kill density per test. The same
pattern appears in the quality metrics: P2 has the lowest compilation and pass
rates of the four conditions, consistent with the model generating more but
less-constrained code. The finding is that \emph{quantity of tests is a poor
proxy for oracle quality}, and it is the strongest single argument in the study
for specifications over generic prompting.

\subsection{Specification Strength and Its Consequences}

Underlying the per-category effect is a distribution of specification
\emph{strength} that the analysis tracks and that explains part of the variation.
A specification is classified as \emph{strong} when both its precondition and its
postcondition are non-trivial, \emph{partial} when exactly one side is non-trivial,
and \emph{weak} when neither is. The classification matters because a strong
specification gives the model two sources of oracle information --- what inputs are
valid and what outputs to expect --- while a partial one gives one and a weak one
gives little more than the signature already did.

The strength distribution varies by category in a way that parallels the
test-generation effect. Accessor and factory methods achieve the highest
proportion of strong specifications, because their behaviour --- read a field,
construct an object --- is fully captured by a postcondition and any input
constraint by a precondition. State-modification methods have the lowest, because
capturing the complete frame condition for a method that mutates object state under
conditional logic is the hardest inference task, and a method whose state change
the engine captures only partially yields a partial rather than a strong
specification. The categories with the strongest specifications are not, however,
the categories with the largest test-generation gains, and the apparent tension
resolves instructively: accessors have strong specifications but small gains
because their signature-only baseline is already high, while control-flow methods
have moderate specification strength but large gains because their signature-only
baseline is low and the specification's guarded postconditions are exactly what the
baseline lacks. The gain is the product of specification strength and baseline
weakness, not of strength alone, which is why the category with the most to gain ---
control flow, with weak baselines and informative guarded postconditions --- shows
the largest improvement despite not having the strongest specifications.

This decomposition refines the headline claim once more. Specifications help in
proportion to the oracle information they add over the baseline, and that
information is the product of how much the specification states (its strength) and
how little the signature already implied (the baseline's weakness). The categories
sort themselves along this product, and the per-category results are its
manifestation.

\subsection{Where Specifications Help Most}

The benefit of specifications varies systematically across method categories.
Control-flow methods benefit most, with a mutation-score improvement of 47.1
percentage points from P1 to P3, because the specification delineates the
expected behaviour of each branch and the model can target each guarded
postcondition with a directed test. Utility methods improve by 40.6 points,
their complex input constraints captured well by inferred preconditions; state-
modification methods improve by 36.6 points, benefiting from postconditions that
state their state changes. Accessors and mutators improve least, by around 19
points, not because specifications fail to help but because their signature-only
baseline is already near the achievable ceiling --- a getter is easy to test
from its signature alone. Within the looping methods, those for which the engine
inferred a non-trivial loop invariant attain a P3 mutation score of roughly 72\%
against roughly 54\% for those left with only a weak invariant, a contrast that
confirms the importance of the loop-invariant heuristics of
Section~\ref{sec:inf-loops} for downstream test quality.

The loop result deserves emphasis because it traces the value of one of the
engine's hardest inference tasks directly to a downstream outcome. A loop invariant
is the single most difficult clause the engine infers, requiring the heuristic
machinery of Section~\ref{sec:inf-loops} where the other clause kinds need only
pattern correspondence, and one might reasonably ask whether the effort is repaid.
The eighteen-point gap between the well-invariant and weak-invariant looping methods
answers that it is: a method whose loop the engine characterises with a non-trivial
invariant yields a specification whose promoted postcondition states what the loop
computes, and a model given that postcondition tests the loop's result directly,
while a method left with only a weak invariant yields a specification that says
little about the loop's outcome and a model that tests it weakly. The downstream
consequence of the loop-invariant heuristics is thus measurable in the mutation
score, and it justifies the disproportionate machinery the engine devotes to loops:
the hardest inference task produces, when it succeeds, the clauses that most
improve the tests, and when it fails produces the methods where P3 most
underperforms its potential. This is the clearest instance in the thesis of an
inference capability's value being traced through to a downstream metric, and it is
the reason loop-invariant strengthening is named among the most consequential
directions for future work.

\section{A Closer Look at Individual Classes}
\label{sec:tg-perclass}

The per-class figures of Table~\ref{tab:tg-results} repay individual attention,
because the classes differ in ways that explain the spread of the effect. The
utility classes \texttt{NumberUtils}, \texttt{BooleanUtils}, and
\texttt{CharUtils} show the largest absolute test counts and large mutation-score
gains from P1 to P3. These classes are dense with conversion and parsing methods
--- turning strings into numbers, characters into integers, booleans into their
string forms --- and such methods have rich input-validity preconditions and
precise output postconditions that the engine captures well. The inferred
specification for a parser states which inputs are valid and what each valid input
yields, and the model turns each clause into a targeted test, which is why these
classes see both many tests and high mutation scores under P3.

The mutable wrapper classes \texttt{MutableInt}, \texttt{MutableDouble}, and
\texttt{MutableBoolean} are interesting for the opposite reason: their methods are
simple state mutations --- increment, add, set --- and their P1 mutation scores
are already moderate because a getter or setter is easy to test from its signature.
Their P3 gains are real but smaller, and concentrated in the methods whose
postconditions relate the new state to the old via \texttt{\textbackslash old}
--- an \texttt{addAndGet} whose contract states the result equals the old value
plus the argument. The \texttt{\textbackslash old}-bearing postconditions are
exactly the clauses a signature cannot imply, and they are where the wrapper
classes' P3 improvement concentrates. \texttt{MutableInt} reaching a 100\%
mutation score under the source condition P4 reflects the simplicity of its
methods: with the body in hand the model tests them exhaustively.

The factory and value classes \texttt{Pair} and \texttt{MutablePair} test the
engine's handling of construction and delegation. Their factory methods carry the
postcondition that the result is non-null and that its components equal the
arguments, and their accessors carry the postcondition that they return the
stored component; the model turns these into construction-and-readback tests that
kill the return-value mutants a signature-only suite would miss. The control-flow
classes \texttt{Validate} and \texttt{Range} are where specifications help most in
relative terms: \texttt{Validate}'s methods are pure guards whose entire behaviour
is a precondition-check-and-throw, so the inferred preconditions \emph{are} the
specification, and \texttt{Range}'s methods are dense with boundary conditions
whose guarded postconditions give the model a test per boundary. These classes'
large P1-to-P3 gains are the per-class manifestation of the category result that
control-flow methods benefit most.

The reading reinforces the central finding from a different angle. Across the
classes, the size of the P3 improvement tracks the gap between what a method's
signature reveals and what its specification states: largest for the parsers and
guards whose behaviour is invisible in their signatures, smallest for the
accessors whose behaviour their signatures already imply. The specification's
value is, consistently, the value of the information it adds over the signature,
and that value varies by class exactly as the information gap does.

\section{Where the Mutants Are Killed}
\label{sec:tg-operators}

The mutation score is an aggregate over a heterogeneous population of seeded
faults, and decomposing it by mutation operator reveals \emph{which kinds} of
fault specifications help catch. The PIT default operator set seeds faults of
several kinds: conditionals-boundary mutations replace a strict comparison with
its non-strict variant; negate-conditionals mutations invert a boolean
condition; math mutations swap a binary arithmetic operator; increments
mutations flip an increment to a decrement; and a family of return-value
mutations replace a method's return with a default, a negated, or a substituted
value. Each kind stresses a different aspect of a test suite's oracle.

The operators that specifications help most are exactly those whose faults change
\emph{what a method returns} rather than \emph{whether it runs}. A return-value
mutation that makes a method return zero where it should return its computed
result survives any test that does not assert the specific expected value --- and
the signature-only and guidance conditions, lacking an oracle, rarely assert the
specific value, so they leave these mutants alive. The specification condition,
carrying postconditions that state the expected result, generates tests that
assert it, and kills the return-value mutants directly. The same holds for
conditionals-boundary and negate-conditionals mutations on methods with guarded
postconditions: the inferred guarded clauses tell the model which input should
produce which result, so a mutation that shifts a branch boundary or inverts a
guard produces an observably wrong result on a test the specification motivated.
This is the operator-level mechanism behind the category result of
Section~\ref{sec:tg-value}: control-flow methods improve most because their
faults are predominantly conditional-boundary and negate-conditional faults,
which are exactly the faults the inferred guarded postconditions equip the model
to catch.

The operators that specifications help least are those whose faults are caught by
mere execution. A void-method-call mutation that removes a call with a
side effect is often caught by any test that exercises the path, because the
missing side effect changes observable state regardless of the oracle's strength;
specifications add little here, because the signature-only baseline already kills
these mutants. The decomposition thus refines the headline result: specifications
do not uniformly raise the mutation score, they raise it specifically on the
return-value and conditional faults that an oracle is required to catch, and they
leave the execution-detectable faults to the coverage that any condition
achieves. The aggregate 40.7-point improvement is concentrated in the
oracle-dependent operators, which is precisely where the oracle problem predicts
it should be.

\section{Discussion}
\label{sec:tg-discussion}

\subsection{What the Two Results Together Establish}

The fidelity result and the value result are independent, and their conjunction
is the chapter's contribution. Fidelity alone would establish that the engine
infers accurate specifications but not that anyone should care; value alone
would establish that specifications help test generation but leave open whether
\emph{inferred} specifications, with their imperfections, help as much as
hand-written ones would. Together they establish the full claim: the engine
infers specifications that are accurate under both a human and a formal oracle,
and supplying those specifications --- imperfections included --- to a language
model improves its test generation by a large and statistically robust margin.
The improvement is achieved with an input an order of magnitude more compact
than the source and is not replicated by generic prompt enrichment, which
isolates the specification as the active ingredient.

\subsection{The Three Tractability Limits}

The 7.1\% of clauses that OpenJML does not discharge are not arbitrary failures;
they fall into three categorically identifiable patterns that exceed the
tractability budget of the solver at the project's strictness configuration.
The first is \emph{recursive multiplicative growth}: the solver's linear-
arithmetic core cannot unfold the inductive product needed to bound a factorial,
a Fibonacci recurrence, or a power function, so the engine emits a literal
parameter bound (a small-input contract) that discharges while the unbounded
form remains unprovable. The second is \emph{quantified invariants over mutable
arrays}: an in-place array scaling loop requires the solver to preserve a
universally quantified element bound across an array store, and three
formulations each exceeded the per-clause time budget. The third is
\emph{the for-each translation}: in a for-each accumulator the loop variable is
bound from the array internally, but the relationship between the variable and
the indexed element is not exposed to user JML, so the per-step overflow check
cannot reference the precondition's element bound. These are properties of the
SMT solver and the JML translation at the present strictness, not of the
engine's heuristics; they delimit what verification technology can currently
certify, not what the engine can emit, and naming them precisely is more useful
than reporting an undifferentiated failure rate.

\subsection{The Tractability Limits, With Examples}

The three tractability limits warrant concrete illustration, because the
distinction between a clause that is false and a clause that is true but
unprovable is central to reading the discharge results correctly. The first limit,
recursive multiplicative growth, is exemplified by a factorial method whose
postcondition would relate the result to the product of the integers up to the
argument. The solver's linear-arithmetic core cannot unfold the inductive product
--- it would need to reason that the result of the recursive call, multiplied by
the argument, satisfies a bound, and the multiplication of a variable by a
variable is outside linear arithmetic. The engine's response is to emit a literal
parameter bound --- the argument must be small enough that the factorial does not
overflow --- which discharges as a small-input contract while the unbounded
relationship remains unprovable. The clause the engine emits is true and provable;
the clause it would like to emit is true but beyond the solver.

The second limit, quantified invariants over mutable arrays, is exemplified by an
in-place array scaling loop that multiplies each element by a factor. The
postcondition would state that every element equals its old value times the
factor, and the loop invariant would state that this holds for the prefix
processed so far. Preserving the universally quantified invariant across the array
store at each iteration requires the solver to reason that writing one element
leaves the bound on the others intact --- a property of the array theory that the
solver supports in principle but that, combined with the quantifier, exceeds the
time budget in practice. Three formulations were attempted, each timing out, which
is the evidence that the limit is the solver's and not a missing formulation the
engine could supply.

The third limit, the for-each translation, is the most purely a verifier artefact.
In a for-each loop over an array, the loop variable is bound internally from the
array at the iteration index, but the relationship between the variable and the
indexed element is not exposed to user-level JML --- the variable simply appears,
disconnected from the array in the specification namespace. A per-step overflow
check on the accumulator therefore cannot reference the precondition's bound on
the array element, because it cannot name the array element the loop variable came
from. This is a gap in the JML translation of the for-each construct, not in the
engine's inference or the solver's power, and it is the case the thesis classifies
as a structural limit of the toolchain rather than of the technique. The three
together account for the bulk of the non-discharged postconditions, and naming
them precisely converts an undifferentiated failure rate into a map of exactly
where the verification frontier lies.

\subsection{Implications}

The results have practical implications beyond the test-generation setting.
That the engine produces useful specifications for 99\% of the corpus at 88
milliseconds per file makes integration into continuous-integration pipelines,
pre-commit hooks, and editor tooling practicable. A concrete workflow makes the
implication tangible. A continuous-integration job runs the engine on every pull
request, infers specifications for the changed methods, discharges them through the
verifier, and reports the flagged clauses as review comments alongside the
diff --- so that a developer reviewing a change sees not only the code but the
contracts the change implies and any clauses that failed to discharge against it. A
clause that fails discharge on a changed method is a signal that the change may have
violated an implicit contract, surfaced at exactly the moment a reviewer is looking
at the code. The 88-millisecond-per-file throughput is what makes this viable: the
inference disappears into the build, and the developer pays no perceptible cost for
the contracts.

The editor-integration scenario is analogous at a finer grain. An editor that runs
the engine as the developer types can display the inferred contract of the method
under the cursor, turning the implicit behaviour the developer is writing into an
explicit contract they can check against their intent --- and flagging the moment
the code and the inferred contract diverge from what the developer meant. These
scenarios are not evaluated in this thesis, which measures the test-generation
consumer specifically, but they follow from the same properties --- accurate
specifications produced fast --- and they indicate that the test-generation result
is one instance of a broader applicability the engine's speed and accuracy enable. The complementarity with
Jdoctor suggests a combined workflow that draws preconditions from both code
structure and documentation. And the central finding --- that compact
behavioural specifications shift a language model's output toward genuine
behavioural checks --- generalises in principle to any model and any JML-aware
consumer, though confirmation across models and codebases remains future work
and is named as a threat below.

A final implication concerns the relationship between this study and the rest of
the thesis. The study measures value to one consumer, a language-model test
generator, chosen because it is the most demanding and most contemporary consumer
of code-level intent. But the value the study measures is not specific to that
consumer; it is the value of the oracle information the specification carries, and
any oracle-dependent consumer --- a search-based test generator given the
specification as assertion oracles, a verifier given it as the contract to check, a
developer given it as documentation to read --- draws on the same information. The
language model is the study's instrument for measuring that value precisely,
because its output can be scored by mutation testing, but the value is a property of
the specification, not of the model. This is why the distribution study of
Chapter~\ref{ch:distribution} matters: the value the model demonstrates is available
to every consumer the distribution mechanism can reach, and the test-generation
result is best read not as a fact about language models but as a measurement, via
language models, of how much intent an inferred specification carries.

\section{Qualitative Analysis of Generated Tests}
\label{sec:tg-qualitative}

The aggregate metrics establish that specifications help; inspecting the
generated tests shows \emph{how}. Consider a representative method, a numeric
parser that returns a default on invalid input:

\begin{lstlisting}
public static int toInt(String str, int defaultValue) {
    if (str == null) { return defaultValue; }
    try { return Integer.parseInt(str); }
    catch (NumberFormatException e) { return defaultValue; }
}
\end{lstlisting}

Under the signature-only condition (P1) the model, lacking any statement of
intended behaviour, generates tests that exercise the obvious happy path and
assert weakly. A typical P1 test passes \texttt{"42"} and asserts the result is
\texttt{42}, and a second passes \texttt{null} and asserts the result equals the
default --- both correct, but both reconstructible from the method name alone,
and neither probing the catch path. The P1 suite covers the method's lines while
leaving its behaviour largely unchecked, which is why its mutation score is low:
a mutant that changed the default-return to return zero would survive every test
that happened to pass zero as the default.

Under the guidance condition (P2) the model, instructed to cover edge cases,
generates many more tests --- empty string, whitespace, leading-plus,
\texttt{Integer.MAX\_VALUE}, overflow strings, embedded letters --- but in the
absence of an oracle it asserts that these ``do not throw'' or returns
``some int'', because it has no statement of what each input \emph{should}
produce. The extra tests inflate the count and the coverage but kill few
additional mutants, because their assertions do not distinguish a correct result
from a plausible wrong one. This is the mechanism behind the headline P2-versus-P3
inversion: P2 generates the most tests and a lower mutation score, because
quantity without an oracle is motion without progress.

Under the specification condition (P3) the model receives the inferred contract
--- the guarded postconditions stating that a null or unparseable string yields
the default and a parseable string yields its integer value --- and generates
tests that assert against those postconditions directly. A P3 test passes
\texttt{"notanumber"} and asserts the result equals the default \emph{value
supplied}, not merely that it does not throw; another passes a valid string and
asserts the exact parsed result. These assertions kill exactly the mutants that
P1 and P2 miss, because they encode what the method should do rather than what it
happens to do. The specification's clause-by-clause structure also shapes the
suite: the model produces roughly one test per guarded postcondition, yielding a
suite that is smaller than P2's but denser in mutant-killing power.

The pattern generalises across categories. For control-flow methods the
guarded postconditions give the model a test target per branch, which is why
control-flow methods show the largest P1-to-P3 improvement. For accessors the
specification adds little the signature did not already imply, which is why their
improvement is smallest. The qualitative inspection thus explains the
category-level quantitative result of Section~\ref{sec:tg-value}: specifications
help most where the gap between what a method does and what its signature reveals
is largest, and least where the signature already determines the behaviour.

\section{Threats to Validity}
\label{sec:tg-threats}

\paragraph{Internal validity.} Sampling non-determinism is mitigated by five
independent runs per method per condition and by statistical aggregation at low
temperature. Prompt confounding is addressed by minimal prompts differing only
in the controlled variation and by the P4 control. Subjectivity in manual
validation is reduced by two independent raters with explicit criteria and high
inter-rater reliability.

\paragraph{External validity.} The corpus is a single, mature, well-structured
library, which may favour the approach, and it lacks the \emph{other} category
of event handlers and callbacks. Results may not transfer to enterprise
applications or domain-specific libraries, and confirmation across heterogeneous
codebases is future work. The model has likely encountered Commons Lang in
training, but the consistent P1-to-P3 gain is inconsistent with simple training-
set recall: were the model merely recalling memorised tests, P1 would already
approach P4, which it does not (27.5\% against 94.8\%). A single model family was
used for budget reasons; the defended contribution is the ordering
P1~$<$~P2~$<$~P3~$<$~P4, which multi-model studies in the literature suggest is
not model-specific, but a full multi-model replication at this scale is left to
future work.

\paragraph{Heuristic soundness.} Inference is heuristic and could in principle
emit malformed or false clauses even after manual review. This is mitigated by
discharging every clause through OpenJML and by the regression suite that gates
engine changes against the verifier; flagged clauses are reviewed rather than
retained. The manual-validation precision and the formal-discharge rate are
reported separately precisely because they are independent oracles catching
different errors.

\paragraph{Construct validity.} Mutation testing has known limitations,
principally equivalent mutants and operator representativeness. The PIT default
operator set is used, equivalent mutants are filtered, and test count is reported
as a complementary metric. The inversion of the count and mutation-score ordering
between P2 and P3 is itself evidence that the mutation score measures something
the count does not.

\paragraph{Conclusion validity.} The principal contrast is well powered: with
312 paired observations and the within-pair variance estimated from a pilot, the
design detects 2.2 points at 80\% power under the Bonferroni-corrected threshold,
and the observed 40.7-point effect yields post-hoc power exceeding 0.999. All
reported significant differences survive the correction across the family of
twenty-four contrasts.

\section{The Generalisation to Service-Boundary Code}
\label{sec:tg-service}

The corpus of this study is pure library code, where the oracle problem is real
but not at its most severe, and it is worth setting out explicitly why the setting
where the problem bites hardest --- service-boundary code --- is both the natural
next target and one for which the present result is a lower bound rather than an
upper one. A method at a service boundary calls a remote endpoint or an RPC stub:
its body delegates to an external service, its signature conveys nothing about the
service's status codes, error semantics, or retry behaviour, and its source, even
when available, reveals only the call, not the contract the remote service
honours. For such a method, the source-code condition P4 --- the upper bound in
this study --- loses its advantage, because the source does not contain the answer;
the answer lives in the remote service's behaviour, which the source does not
describe.

This changes the structure of the argument in a way that favours specifications. In
the present study, P4 dominates P3 because the source contains the test answer
outright; at a service boundary, the source does not, so the gap between P3 and P4
should narrow or close, and the inferred or documented specification of the
service's contract --- carried, as Chapter~\ref{ch:distribution} shows, in the
service's OpenAPI description --- becomes the only source of oracle information
that signature, source, and documentation do not already encode. The hypothesis
for the planned follow-up is therefore that the P1-to-P3 gain reported here will
\emph{compound} at the service boundary: where the source is uninformative about
intended behaviour, the specification is not merely a more compact alternative to
the source but the sole carrier of the oracle, and signature-only generation
collapses into mocking the call with no behavioural check at all. The present study
establishes the effect in the setting where the source competes with the
specification; the follow-up will test it in the setting where the source cannot.
